% !PRINTED 345
\weberpage{345}

%<<CONT
worin $\varphi, \psi$ ganze ganzzahlige Functionen der Variablen $\xi, \eta$
sind, deren Grad in Bezug auf keine der Variablen den Werth 2
übersteigt. Ordnen wir diese Functionen nach $\eta$, so können
wir setzen:
\weq{5}{\xi' = A\eta^2 + B\eta + C, \quad \eta' = A'\eta^2 + B'\eta + C',}
worin $A, B, C, A', B', C'$ ganze Functionen von $\xi$, höchstens
vom $2^{\text{ten}}$ Grade sind. Wir wollen in den Gleichungen (4) die
Variable $\xi$ festhalten und $\eta = 0, 1, 2$ setzen, d.\,h. wir wenden
(5) auf das Tripel $(\xi, 0)\,(\xi, 1)\,(\xi, 2)$ an. Die entsprechenden
$(\xi_1', \eta_1'), (\xi_2', \eta_2'), (\xi_3', \eta_3')$ müssen dann den Relationen (3)
genügen, also:
\[
\begin{array}{ll}
\xi_1' \equiv C, & \eta_1' \equiv C',\\
\xi_2' \equiv A + B + C, & \eta_2' \equiv A' + B' + C',\\
\xi_3' \equiv A - B + C, & \eta_3' \equiv A' - B' + C',
\end{array}
\]
und daraus folgt, dass $A$ und $A'$ für jedes $\xi$ congruent mit
Null sein müssen.

Ebenso kann man schliessen, dass in den Substitutionen (4)
kein Glied mit $\xi^2$ vorkommen kann, und dass daher diese Substitutionen
 die Form haben müssen:
\[
\begin{array}{l}
\xi' \equiv m\,\xi\,\eta + a\,\xi + b\,\eta + c\\
\eta' \equiv m'\,\xi\,\eta + a'\,\xi + b'\,\eta + c'
\end{array} \quad (\mathrm{mod}\ 3).
\]
Wenn man aber diese Substitutionen auf das Tripel $(0,\,0)\,(1,\,1)\,(2,\,2)$
anwendet, und die Summe der entsprechenden $\xi'$ und $\eta'$ gleich
Null setzt, so folgt, dass $m = m' = 0$ sein muss, und dass also
jede Permutation der Gruppe $P$ in der Form
\weq{6}{\begin{array}{l}
\xi' \equiv a\,\xi + b\,\eta + c\\
\eta' \equiv a'\,\xi + b'\,\eta + c'
\end{array} \quad (\mathrm{mod}\ 3)}
enthalten sein muss, wie behauptet war.

\gesperrt{2. Umgekehrt gehen durch jede lineare Substitution
von der Form (6) drei Grössen $(\xi, \eta)$ eines Tripels
in drei $(\xi', \eta')$ über, die gleichfalls ein Tripel
bilden.}

Denn aus (6) folgt für irgend drei Zahlenpaare $(\xi_1, \eta_1)$,
$(\xi_2, \eta_2), (\xi_3, \eta_3)$:
\weq{7}{\begin{array}{l}
\xi_1' + \xi_2' + \xi_3' \equiv a\,(\xi_1 + \xi_2 + \xi_3) + b\,(\eta_1 + \eta_2 + \eta_3)\\
\eta_1' + \eta_2' + \eta_3' \equiv a'\,(\xi_1 + \xi_2 + \xi_3) + b'\,(\eta_1 + \eta_2 + \eta_3)
\end{array} \quad (\mathrm{mod}\ 3);}
%CONT>>
